\section{The ATLAS method}\label{sec:method}

\subsection{E-processes and Ville's inequality}

An \emph{e-variable} for a null is a nonnegative random variable with expectation at most one under
it. A \emph{test supermartingale} is a process $W_0=1$, $W_k=\prod_{j\le k}e_j$ with $\E[e_j\mid\Fcal_{j-1}]\le1$
under the null. The engine of time-uniform validity is classical.

\begin{lemma}[Ville \citeyear{ville1939}]\label{lem:ville}
If $(W_k)_{k\ge0}$ is a nonnegative supermartingale with $W_0=1$ then $\PR(\sup_kW_k\ge1/\alpha)\le\alpha$.
\end{lemma}

\subsection{The per-horizon betting e-process}\label{sec:betting}

Fix a horizon $h$ and a statistic $Z_k\in[z_{\mathrm{lo}},z_{\mathrm{hi}}]$ calibrated at $\eps_h$
(Definition~\ref{def:cert}), so that the faithfulness null implies
\begin{equation}\label{eq:null-mean}
\E[Z_k\mid\Fcal_{k-1}]\le\eps_h\qquad\text{for all }k.
\end{equation}
For a predictable bet $\lambda_k\in[0,\lambda_{\max}]$ with $\lambda_{\max}=1/(z_{\mathrm{hi}}-z_{\mathrm{lo}})$ set
\begin{equation}\label{eq:evar}
e_k\defeq1+\lambda_k(Z_k-\eps_h),\qquad W^{(h)}_K\defeq\prod_{k\le K}e_k .
\end{equation}

\begin{proposition}[Per-round validity]\label{prop:evar}
If \eqref{eq:null-mean} holds and $\eps_h\in[z_{\mathrm{lo}},z_{\mathrm{hi}}]$ then $e_k\ge0$ and
$\E[e_k\mid\Fcal_{k-1}]\le1$; hence $W^{(h)}$ is a test supermartingale.
\end{proposition}
\begin{proof}
$e_k\ge1-\lambda_{\max}(\eps_h-z_{\mathrm{lo}})\ge0$, and
$\E[e_k\mid\Fcal_{k-1}]=1+\lambda_k(\E[Z_k\mid\Fcal_{k-1}]-\eps_h)\le1$ because $\lambda_k$ is predictable.
\end{proof}

The bet governs power, not validity. We use a uniform mixture over a grid
$\lambda^{(1)},\dots,\lambda^{(m)}\in(0,\lambda_{\max}]$: a convex combination of test supermartingales
is a test supermartingale, so $W^{(h)}_K=\frac1m\sum_j\prod_{k\le K}(1+\lambda^{(j)}(Z_k-\eps_h))$ is valid, parameter-free,
and adapts to an unknown margin at a cost of $\log m$ in the log-wealth. Section~\ref{sec:exp-abl} compares it with a fixed bet and a
plug-in bet.

\subsection{Multi-step validity: why chains, and why they cannot be pooled}\label{sec:overlap}

If a horizon-$h$ prediction is formed at \emph{every} round, consecutive residuals share the innovations in
a window of length $h-1$, so the residual stream is moving-average of order $h-1$. Conditioning on the past
reveals part of the current residual and \eqref{eq:null-mean} fails in general even for a faithful
model. The same effect underlies the overlap problem in sequential multi-step forecast evaluation
\citep{henzi2022correction}. We restore \eqref{eq:null-mean} by running each horizon on
\emph{non-overlapping chains} of issue times $c,c+h,c+2h,\dots$: the $k$-th chain element is resolved before the next
is issued, so $\Fcal_{k-1}$ contains everything the next residual could depend on. The cost is real: horizon $h$ receives a factor $h$ fewer
updates.

One might try to win these updates back by running all $h$ offset chains and rejecting when
any crosses its threshold (a union bound, costing $\log h$). That is valid but it does not pool
evidence: each chain sees only a $1/h$ share of the data and its wealth grows at the same rate per update
as the single chain, so taking the first of $h$ crossings buys at most a small order-statistic
gain that the $\log h$ in the threshold cancels. Averaging the chains' wealth would pool, but
chains live on different filtrations, so the average is an e-value at fixed times and not an e-process (an
optional-stopping argument does not go through). Section~\ref{sec:exp-horizon} measures the
interleaved union: it is valid and no faster than the single chain. We therefore keep the simple estimator
and treat the factor $h$ as the intrinsic price of certifying long horizons directly.

\subsection{Neural world models: sampled predictives and committed action blocks}\label{sec:neural}

A neural world model rarely has a closed-form multi-step predictive. We take it by sampling: the stochastic model
(for instance a Gaussian head on the next latent) is rolled forward $n$ times under the action sequence, and the standardized excess of
Example~\ref{ex:gauss} is computed from the sample mean and covariance on its non-degenerate subspace. Two requirements for validity
follow from Section~\ref{sec:setup}. First, for $h>1$ the action sequence must be fixed at issue time (for $h=1$ the action is already predictable). A planner that replans at every step reacts
to intermediate noise, so the realized actions are not independent of the outcome and the null can fail even for a faithful model.
We therefore let the planner \emph{commit} to an open-loop block of length $L$ and replan only at block boundaries; for
each horizon $h\le L$ the block yields $\lfloor L/h\rfloor$ non-overlapping chain elements, issued at block offsets $0,h,2h,\dots$,
each conditioned on actions that were known when it was issued. Second, the statistic's conditional mean is tied to the divergence only if the sampled predictive's covariance is
adequate, which is why the tolerance is calibrated on held-out nominal closed-loop episodes of the same controller (Section~\ref{sec:exp-control}). A tolerance estimated from data is not covered by Theorem~\ref{thm:validity}: with it, validity is empirical, and the guarantee is conditional on the estimate being an upper bound on the statistic's conditional mean.

The commit length is the natural place for the gate: the planner runs with $L=\min(L_{\max},\hstar(t))$, checks the frontier after every executed step, and
aborts the remaining block and enters a safe fallback when $\hstar\le1$ (Algorithm~\ref{alg:atlas}).

\subsection{Two outputs: the revocation frontier and the certified horizon}\label{sec:frontier}

\paragraph{Revocation frontier (monitor).} Run one e-process per horizon at level $\alpha/H$ and revoke horizon $h$ at
$\tau_h=\inf\{k:W^{(h)}_k\ge H/\alpha\}$, permanently (Ville concerns the event of \emph{ever} crossing). Define
\begin{equation}\label{eq:hstar}
\hstar(t)\defeq\max\{h:\ \text{no horizon }h'\le h\text{ is revoked by round }t\},
\end{equation}
with $\hstar=0$ if horizon one is revoked. The prefix rule encodes that a broken short horizon casts
doubt on longer ones, and matches the use as a depth clip. The guarantee (Theorem~\ref{thm:validity}) is
\emph{no false revocation}: a faithful horizon is lost with probability at most $\alpha$.

\paragraph{Certified horizon (certificate).} In estimation mode the same machinery gives a
time-uniform upper bound $U^{(h)}_t$ on the running average of conditional means of horizon $h$ (Lemma~\ref{lem:cs}),
hence, for a certifiable statistic, an upper bound on the average divergence and, through the modulus, on total
variation: $\widehat{\TV}_h(t)\defeq\min\{1,\phi(cU^{(h)}_t+\bar\rho)\}$. For a budget $\tau$ the certified horizon is
\begin{equation}\label{eq:hcert}
\hcert_\tau(t)\defeq\max\{h:\ \widehat{\TV}_i(t)\le\tau\ \text{for all }i\le h\}.
\end{equation}
The guarantee (Theorem~\ref{thm:cert}) is the opposite error: with probability at least $1-\alpha$, simultaneously
over rounds and horizons, the true average discrepancy of every horizon inside $\hcert_\tau(t)$ is at most $\tau$.
The planner's gate is $\min(\hstar(t),\hcert_\tau(t))$. The two do not substitute for each other: $\hstar$ can
be high on a model that is mediocre but stable, and $\hcert_\tau$ can be low on a faithful model that has not yet
accumulated enough evidence.

\paragraph{Change detection.} To detect a change in faithfulness rather than test a fixed null we use the
Shiryaev--Roberts e-detector \citep{shin2024edetectors}, $R_t=(1+R_{t-1})e_t$ (with the mixture bet), which accumulates
evidence over all candidate changepoints so that wealth lost before the change cannot delay the alarm. Algorithm~\ref{alg:atlas} summarizes the procedure.

\begin{algorithm}[t]
\caption{ATLAS: monitoring and certification of a deployed world model}\label{alg:atlas}
\begin{algorithmic}[1]
\State \textbf{input:} horizons $\{1,\dots,H\}$, tolerances $\{\eps_h\}$, level $\alpha$, bet grid $\{\lambda^{(j)}\}$, budget $\tau$, constants $(c,\bar\rho)$
\State initialize wealth $W^{(h)}\gets1$, $\texttt{revoked}(h)\gets\text{false}$, and a confidence sequence $\mathrm{CS}^{(h)}$ for each $h$
\For{each round $t=1,2,\dots$}
  \For{each horizon $h$ whose chain element resolves at $t$}
    \State form the bounded statistic $Z$; update $W^{(h)}\gets\sum_j\tfrac1m\Pi_j(1+\lambda^{(j)}(Z-\eps_h))$ and $\mathrm{CS}^{(h)}$
    \If{$W^{(h)}\ge H/\alpha$} \State $\texttt{revoked}(h)\gets\text{true}$ \EndIf
  \EndFor
  \State $\hstar(t)\gets$ largest prefix with no revoked horizon; $\hcert_\tau(t)\gets$ largest prefix with $\widehat{\TV}_i\le\tau$
  \State \textbf{output} the gate $\min(\hstar(t),\hcert_\tau(t))$
\EndFor
\end{algorithmic}
\end{algorithm}
